% !TeX root = ../main.tex

Die Ergebnisse zeigen, dass der Retry-Ansatz grundsätzlich geeignet ist, um die Erfolgsrate von den beschriebenen
Fehlerklassen zu erhöhen.
Der Nutzen eines Retries unterscheidet sich jedoch deutlich zwischen den Fehlerklassen.
Gleichzeitig sorgt der Retry für zusätzliche Ausführungskosten, sodass eine möglichst hohe Erfolgsrate nicht
unabhängig von der dafür benötigten Ausführungszeit betrachtet werden kann.
Eine allgemeingültige Konfiguration für Retries ist daher nicht ableitbar.
Stattdessen muss die Strategie und ihre Parameter sorgfältig an die jeweilige Fehlerklasse angepasst werden.
Darüber hinaus zeigen die Ergebnisse einen negativen Einfluss steigender Concurrency auf die Erfolgsraten.
Unter den untersuchten Bedingungen ist insbesondere ab einer Concurrency von 5 ein deutlicher Rückgang der
Erfolgsraten zu beobachten.
Dies deutet darauf hin, dass der Retry-Ansatz bei zunehmender Konkurrenz um Datenbankressourcen an seine Grenzen
gelangt.

Ohne Retries führt jeder Konflikt zum Abbruch.
Da bei höherer Concurrency die Wahrscheinlichkeit für Kollision steigt, fällt die Erfolgsrate bei \texttt{s0} stark ab.
Die Strategien mit Retries puffern die Last ab, geraten aber mit steigender Concurrency an die Ressourcengrenze.
Bei niedriger Concurrency profitieren Deadlocks stark von \texttt{s1}, da der Konflikt nach dem Abbruch meistens
sofort aufgelöst ist.
Erwähnenswert ist die gleiche Erfolgsrate von Deadlocks und Serialization Failures bei \texttt{cc2} von 50.08\,\%,
die darauf zurückzuführen ist, dass bei zwei Clients genau einer gewinnt.
Unter hoher Concurrency brechen jedoch alle Strategien ein, da neue Deadlocks entstehen können.
Für Serialization Failures sind \texttt{s2} und \texttt{s3} robuster bei höherer Concurrency ab \texttt{cc10}.
Ein Delay kann also die zeitgleiche Konkurrenz abmildern.
Die Lock Timeouts verschärfen sich durch hohe Concurrency dramatisch.
Wartende Transaktionen blockieren Ressourcen, wodurch auch eine Strategie mit Delay bei \texttt{cc20} kaum noch Gewinn
bringt.
Concurrency verstärkt den Unterschied zwischen den Strategien.
Strategien mit Delay skalieren unter Last besser für Serialization Failures.
Bei Deadlocks hingegen ist der \emph{immediate Retry} bei geringer Last den anderen Strategien etwas überlegen.
Um zu verstehen, warum die einzelnen Strategien unter variierender Last dieses spezifische Verhalten zeigen,
müssen deren Funktionsweisen für die jeweiligen Fehlerklassen genauer betrachtet werden.

Wie in Kapitel~\ref{subsec:s} beschrieben bringen unterschiedliche Strategien unterschiedliche Erfolgsraten für die
einzelnen Fehlerklassen.
Für Deadlocks ist Strategie 1 die mit der höchsten Rate.
Eine mögliche Erklärung für das gute Abschneiden des Immediate Retry ist, dass sich der Konflikt nach dem Auflösen
des Deadlocks gelöst hat.
Der zusätzliche Delay könnte eher wieder dafür sorgen, dass eine Transaktion in einen anderen Deadlock verwickelt wird.
Dies könnte erklären, warum die Erfolgsrate für die \texttt{s2} und \texttt{s3} wieder sinkt.
Serialization Failures haben die höchste Erfolgsrate hingegen bei \texttt{s2}.
Das bessere Abschneiden eines Retries mit Delay deutet darauf hin, dass ein direkter Wiederholungsversuch bei der
Fehlerklasse häufiger erneut mit einer Transaktion konkurriert.
Ein kurzer Delay kann der konkurrierenden Transaktion zusätzliche Zeit geben ihren Vorgang abzuschließen.
Bei Lock Timeouts hat auch \texttt{s2} die beste Erfolgsrate.
Da die Transaktion vorher auf ein Lock gewartet hat, kann ein Delay dabei helfen nicht wieder in diese Situation zu
geraten.
Der starke Anstieg der Ausführungszeiten nicht erfolgreicher Transaktionen steht dabei mit den zusätzlichen
Wiederholungsversuchen in Zusammenhang.
Die höheren Ausführungszeiten erfolgreicher Transaktionen lassen sich durch die in der Testumgebung künstlich
verlängerte Transaktionsdauer erklären.
Dies könnte auch die niedrigen Erfolgsraten erklären.

Wie bereits ausgeführt, lässt sich eine gute Strategie nicht nur über die Erfolgsrate definieren.
Es wäre beispielsweise sinnlos eine Erfolgsrate von 90\,\% zu erreichen, wenn eine Transaktion dafür 1000 Wiederholungen
durchführen muss und dadurch zwei Stunden benötigt.
Es muss also ein Punkt gefunden werden, bei der die Erfolgsrate steigt, die Ausführungszeit sich aber noch in einem
möglichst vertretbaren Rahmen befindet.

% !TeX root = ../main.tex

\begin{table}[htbp]
	\centering
	\smallsize
	\begin{tabular}{l|c|ccc}
		\toprule
		& & \multicolumn{3}{c}{Faktoren} \\
		\cmidrule(lr){3-5}
		Type & S & Erfolgsrate-Faktor & Ausführungszeit-Faktor & Effizienzverhältnis \\
		\midrule
		\multirow{3}{*}{Deadlock}
		& s1 & 2.89 & 2.38 & 1.21 \\
		& s2 & 2.61 & 2.24 & 1.17 \\
		& s3 & 2.56 & 2.31 & 1.08 \\
		\midrule
		\multirow{3}{*}{Serialization}
		& s1 & 2.32 & 3.08 & 0.75 \\
		& s2 & 2.44 & 3.45 & 0.71 \\
		& s3 & 2.20 & 2.60 & 0.85 \\
		\midrule
		\multirow{3}{*}{Lock Timeouts}
		& s1 & 2.60 & 1.09 & 2.39 \\
		& s2 & 3.00 & 1.20 & 2.50 \\
		& s3 & 2.80 & 1.19 & 2.35 \\
		\bottomrule
	\end{tabular}
	\caption{Erfolgsrate und Ausführungszeit relativ zu s0 sowie deren Verhältnis.}
	\label{tab:rDeadlocks}
\end{table}

Die Tabelle~\ref{tab:rDeadlocks} zeigt beispielsweise das Verhältnis von Erfolgsrate zur Ausführungszeit für die
einzelnen Strategien gruppiert nach den Fehlerklassen.
Die Strategien werden anhand des Verhältnisses zwischen dem Gewinn an Erfolgsrate und dem damit verbundenen Anstieg
der Ausführungszeit bewertet.
Dabei erfolgt der Vergleich jeweils innerhalb einer Fehlerklasse, da sich die absoluten Ausführungszeiten der
verschiedenen Fehlerklassen aufgrund ihrer unterschiedlichen Eigenschaften nicht direkt miteinander
vergleichen lassen.
Bei allen Fehlerklassen führt bereits ein möglicher Retry zu einer deutlichen Steigerung der Erfolgsrate.
Im Verhältnis zur zusätzlich benötigten Ausführungszeit sind die Strategien jedoch unterschiedlich effizient.
So ist die beste Strategie für Deadlocks die Strategie1, so wie es die Erfolgsrate schon zeigt.
Bei Lock Timeouts hat die Strategie mit der höchsten Erfolgsrate ebenfalls das beste Verhältnis zur Ausführungszeit.
Serialization Failures hingegen haben bei der höchsten Erfolgsrate nicht das beste Verhältnis.
Generell erreicht bei dieser Fehlerklasse keine Strategie ein Verhältnis von größer 1.
Gemessen daran bietet somit keine der untersuchten Retry-Strategien gegenüber \texttt{s0} einen Vorteil.
Wenn sich jedoch auf eine Retry-Strategie festgelegt werden sollte, so wäre für Serialization Failures die dritte
Strategie noch die beste Strategie, da sie das höchste Verhältnis hat.

Anhand dieses Konzepts lassen sich die besten Parameter für die einzelnen Strategien herausfinden.
In Tabelle~\ref{tab:bestConfig} wird dies dargestellt.
Für Serialization Failures bietet ein Retry gemäß dem gewählten Bewertungsverhältnis keinen Vorteil gegenüber
\texttt{s0}.
Für die weitere Betrachtung wird dennoch \texttt{s3} als beste Retry-Strategie herangezogen, da sie innerhalb der
Strategien mit Wiederholungsversuch das beste Verhältnis aus Erfolgsrate und Ausführungszeit erreicht.

% !TeX root = ../main.tex

\begin{table}[htbp]
	\centering
	\smallsize
	\begin{tabular}{l|lcc}
		\toprule
		Type & Strategy & Delay & Retry-Count \\
		\midrule
		Deadlocks     & s1 & -- & 1 \\
		Serialization & s3 & 0.005\,s & 1 \\
		Lock Timeout  & s2 & 0.1\,s & 5 \\
		\bottomrule
	\end{tabular}
	\caption{Beste Konfigurationen für die Fehlerklassen bei Gleichgewichtung.}
	\label{tab:bestConfig}
\end{table}
Aus der Tabelle sind die theoretisch besten Konfigurationen unter den gegebenen Bedingungen abzulesen.
Dazu muss bei dieser Bewertungsmethode beachtet werden, dass bei niedrigen Werten Änderungen großen Einfluss
auf die relativen Werte haben.
Dadurch kann ein Anstieg von 32\,\% der Erfolgsrate für 110\,ms Mehraufwand erstrebenswert sein, obwohl das Verhältnis
dagegen spricht.
Die Tabellen~\ref{appsec:8.1} zeigen die genauen Werte, aus denen diese Parameter gewählt wurden.
Zusätzlich wurden die Parameter einzeln betrachtet, sodass der Einfluss dieser aufeinander nicht berücksichtigt wurde.
So ist für Deadlocks die beste Konfiguration \texttt{s1} mit einer möglichen Wiederholung.
Für Serialization Failures ist die beste Strategie mit Wiederholen einer Transaktion \texttt{s3} mit 0.005\,s Delay
und ebenfalls einem möglichen Retry.
Da eine Retry-Strategie bei dieser Fehlerart im Verhältnis keinen Nutzen hat ist es eine Überlegung wert, keinen
Retry zu verwenden.
Wenn jedoch eine Retry-Strategie verwendet werden soll, so wäre die in der Tabelle stehende Konfiguration die Beste.
Bei Lock Timeouts ist \texttt{s2} die beste Strategie und die Parameter sollten mit 0.1\,s Delay und fünf möglichen
Retries gewählt werden.
Diese Konfigurationen setzen dabei voraus, dass Erfolgsrate und Ausführungszeit gleich gewichtet werden.
Wird einer der beiden Metriken eine höhere Bedeutung gegeben, kann sich die bevorzugte Konfiguration
entsprechend verschieben.

In der Problemanalyse wurden Anforderungen an einen Retry-Mechanismus abgeleitet, die anhand der Ergebnisse
abschließend bewertet werden können.
Die Anforderung an die Zuverlässigkeit wird durch die Ergebnisse grundsätzlich bestätigt.
Die untersuchten Retry-Strategien können die Erfolgsrate nach einem Fehler deutlich erhöhen.
Gleichzeitig zeigen die Ergebnisse, dass ein Retry nicht für jede Fehlerklasse gleichermaßen geeignet ist.
Insbesondere bei Serialization Failures bietet ein Retry unter dem gewählten Bewertungsmaßstab keinen Vorteil
gegenüber einer Ausführung ohne Wiederholungsversuch.
Auch die Anforderung an die Performance wird durch die Ergebnisse bestätigt.
Mit zunehmender Anzahl möglicher Retries steigen in der Regel die Ausführungskosten, während der zusätzliche Gewinn
der Erfolgsrate ab einem bestimmten Punkt nur noch geringer ausfällt.
Ebenso zeigt sich, dass der Delay einen Einfluss auf das Verhältnis zwischen Erfolgsrate und Ausführungszeit hat.
Retry-Count und Delay stellen damit wesentliche Parameter dar, die kontrolliert und an die jeweilige Fehlerklasse
angepasst werden müssen.
Die Anforderung an die Flexibilität wird durch die Ergebnisse unterstützt.
Diese zeigen, dass unterschiedliche Fehlerklassen von unterschiedlichen Strategien profitieren.
Während für Deadlocks ein unmittelbarer Wiederholungsversuch das beste Verhältnis zwischen Erfolgsrate und
Ausführungszeit erreicht, ist für Lock Timeouts und unter den untersuchten Bedingungen auch für Serialization
Failures eine Strategie mit Delay vorteilhafter beziehungsweise die beste verfügbare Retry-Strategie.

Insgesamt bestätigen die Ergebnisse damit die in der Problemanalyse formulierte Notwendigkeit, Fehlertoleranz und
Systembelastung gemeinsam zu betrachten.
Eine maximale Erfolgsrate stellt nicht automatisch die beste Konfiguration dar.
Stattdessen muss ein angemessenes Verhältnis zwischen zusätzlicher Erfolgswahrscheinlichkeit und den dadurch
entstehenden Ausführungskosten gefunden werden.
Die Evaluation zeigt somit, dass ein konfigurierbarer und fehlerklassenspezifischer Retry-Ansatz grundsätzlich
geeignet ist, diese Anforderungen zu erfüllen.
